\textbf{Phần 1. Trắc nghiệm lựa chọn}
\setcounterpageref{ex}{0}

\begin{ex} %[5D10V62]
Một ô tô khởi hành từ A lúc $7\text{ giờ } 30\text{ phút}$ để đi đến B. Sau khi đi được $1$ giờ $45$ phút thì xe dừng nghỉ $15$ phút, sau đó đi tiếp $2$ giờ $10$ phút nữa thì tới B. Ô tô tới B vào thời điểm nào?
\choice
{$11\text{ giờ } 30\text{ phút}$}
{\True $11\text{ giờ } 40\text{ phút}$}
{$11\text{ giờ } 45\text{ phút}$}
{$12\text{ giờ}$}
\loigiai{
Tổng thời gian đi và nghỉ là:
$$1\text{ giờ } 45\text{ phút} + 15\text{ phút} + 2\text{ giờ } 10\text{ phút} = 4\text{ giờ } 10\text{ phút}.$$
Thời điểm ô tô tới B là:
$$7\text{ giờ } 30\text{ phút} + 4\text{ giờ } 10\text{ phút} = 11\text{ giờ } 40\text{ phút}.$$
\textbf{Đáp án của câu hỏi là B}
}
\end{ex}

\begin{ex} %[5D10V62]
Một người đi xe máy từ A đến B với vận tốc $36\text{ km/h}$. Khi từ B trở về A, người đó tăng vận tốc lên $45\text{ km/h}$ nên thời gian về ít hơn thời gian đi $20$ phút. Độ dài quãng đường AB là
\choice
{$50\text{ km}$}
{\True $60\text{ km}$}
{$72\text{ km}$}
{$90\text{ km}$}
\loigiai{
Tỉ số vận tốc lúc đi và lúc về là:
$$\dfrac{36}{45} = \dfrac{4}{5}.$$
Vì trên cùng quãng đường AB nên thời gian tỉ lệ nghịch với vận tốc, tỉ số thời gian lúc đi và lúc về là $\dfrac{5}{4}$.\\
Hiệu thời gian giữa lúc đi và lúc về là $20\text{ phút} = \dfrac{1}{3}\text{ giờ}$.\\
Thời gian lúc đi là:
$$20 \times 5 = 100\text{ (phút)} = \dfrac{5}{3}\text{ (giờ)}.$$
Độ dài quãng đường AB là:
$$36 \times \dfrac{5}{3} = 60\text{ (km)}.$$
\textbf{Đáp án của câu hỏi là B}
}
\end{ex}

\begin{ex} %[5D10V62]
Hai xe máy cùng xuất phát từ hai tỉnh A và B cách nhau $165\text{ km}$ và đi ngược chiều nhau. Xe thứ nhất đi với vận tốc $45\text{ km/h}$, xe thứ hai đi với vận tốc bằng $\dfrac{2}{3}$ vận tốc xe thứ nhất. Sau bao lâu hai xe gặp nhau?
\choice
{$2\text{ giờ}$}
{\True $2\text{ giờ } 12\text{ phút}$}
{$2\text{ giờ } 15\text{ phút}$}
{$2\text{ giờ } 30\text{ phút}$}
\loigiai{
Vận tốc của xe thứ hai là:
$$45 \times \dfrac{2}{3} = 30\text{ (km/h)}.$$
Tổng vận tốc của hai xe là:
$$45 + 30 = 75\text{ (km/h)}.$$
Thời gian hai xe gặp nhau là:
$$165 : 75 = 2{,}2\text{ (giờ)} = 2\text{ giờ } 12\text{ phút}.$$
\textbf{Đáp án của câu hỏi là B}
}
\end{ex}

\begin{ex} %[5D10V62]
Một ca nô đi xuôi dòng từ bến A đến bến B hết $1$ giờ $20$ phút. Biết quãng sông AB dài $32\text{ km}$ và vận tốc dòng nước là $3\text{ km/h}$. Vận tốc của ca nô khi đi ngược dòng từ B về A là
\choice
{$16\text{ km/h}$}
{\True $18\text{ km/h}$}
{$20\text{ km/h}$}
{$21\text{ km/h}$}
\loigiai{
Đổi: $1\text{ giờ } 20\text{ phút} = \dfrac{4}{3}\text{ giờ}$.\\
Vận tốc xuôi dòng là:
$$32 : \dfrac{4}{3} = 24\text{ (km/h)}.$$
Vận tốc thực của ca nô là:
$$24 - 3 = 21\text{ (km/h)}.$$
Vận tốc ca nô khi ngược dòng là:
$$21 - 3 = 18\text{ (km/h)}.$$
\textbf{Đáp án của câu hỏi là B}
}
\end{ex}

\begin{ex} %[5D10V62]
Một xe tải đi từ A đến B với vận tốc $40\text{ km/h}$. Sau $1$ giờ $30$ phút, một xe con cũng đi từ A đuổi theo xe tải với vận tốc $60\text{ km/h}$. Chỗ xe con đuổi kịp xe tải cách điểm xuất phát A là
\choice
{$150\text{ km}$}
{\True $180\text{ km}$}
{$200\text{ km}$}
{$240\text{ km}$}
\loigiai{
Đổi: $1\text{ giờ } 30\text{ phút} = 1{,}5\text{ giờ}$.\\
Khi xe con bắt đầu đi, xe tải đã đi được quãng đường là:
$$40 \times 1{,}5 = 60\text{ (km)}.$$
Mỗi giờ xe con rút ngắn khoảng cách với xe tải:
$$60 - 40 = 20\text{ (km)}.$$
Thời gian để xe con đuổi kịp xe tải là:
$$60 : 20 = 3\text{ (giờ)}.$$
Quãng đường từ A đến điểm gặp nhau là:
$$60 \times 3 = 180\text{ (km)}.$$
\textbf{Đáp án của câu hỏi là B}
}
\end{ex}

\begin{ex} %[5D10V62]
Một người đi từ A đến B gồm một đoạn đường bằng và một đoạn đường dốc. Người đó đi đoạn đường bằng mất $40$ phút với vận tốc $12\text{ km/h}$ và đi đoạn đường dốc mất $30$ phút với vận tốc $6\text{ km/h}$. Vận tốc trung bình của người đó trên cả quãng đường AB là
\choice
{$9\text{ km/h}$}
{\True $9{,}43\text{ km/h}$}
{$9{,}5\text{ km/h}$}
{$10\text{ km/h}$}
\loigiai{
Đổi: $40\text{ phút} = \dfrac{2}{3}\text{ giờ}$; $30\text{ phút} = 0{,}5\text{ giờ}$.\\
Quãng đường đoạn bằng: $12 \times \dfrac{2}{3} = 8\text{ (km)}$.\\
Quãng đường đoạn dốc: $6 \times 0{,}5 = 3\text{ (km)}$.\\
Tổng quãng đường AB: $8 + 3 = 11\text{ (km)}$.\\
Tổng thời gian đi: $40 + 30 = 70\text{ phút} = \dfrac{7}{6}\text{ giờ}$.\\
Vận tốc trung bình: $11 : \dfrac{7}{6} = \dfrac{66}{7} \approx 9{,}43\text{ (km/h)}$.\\
\textbf{Đáp án của câu hỏi là B}
}
\end{ex}

\begin{ex} %[5D10V62]
Một đoàn tàu hỏa chạy qua một cây cột điện bên đường mất $12$ giây. Cũng với vận tốc đó, đoàn tàu chạy qua một cây cầu dài $360\text{ m}$ mất $30$ giây. Chiều dài của đoàn tàu hỏa là
\choice
{$200\text{ m}$}
{\True $240\text{ m}$}
{$280\text{ m}$}
{$180\text{ m}$}
\loigiai{
Thời gian đoàn tàu đi hết chiều dài cây cầu $360\text{ m}$ là:
$$30 - 12 = 18\text{ (giây)}.$$
Vận tốc của đoàn tàu là:
$$360 : 18 = 20\text{ (m/s)}.$$
Chiều dài của đoàn tàu là:
$$20 \times 12 = 240\text{ (m)}.$$
\textbf{Đáp án của câu hỏi là B}
}
\end{ex}

\begin{ex} %[5D10V62]
Một người đi bộ vòng quanh một quảng trường hình vuông có chu vi $1{,}2\text{ km}$. Người đó đi hết một vòng trong thời gian $18$ phút. Vận tốc của người đó là
\choice
{$3{,}6\text{ km/h}$}
{\True $4\text{ km/h}$}
{$4{,}2\text{ km/h}$}
{$4{,}5\text{ km/h}$}
\loigiai{
Đổi: $18\text{ phút} = 0{,}3\text{ giờ}$.\\
Vận tốc của người đó là:
$$v = 1{,}2 : 0{,}3 = 4\text{ (km/h)}.$$
\textbf{Đáp án của câu hỏi là B}
}
\end{ex}

\begin{ex} %[5D10V62]
Lúc $6\text{ giờ } 15\text{ phút}$, một người đi xe đạp từ nhà ra ga tàu với vận tốc $12\text{ km/h}$. Người đó đến ga lúc $7\text{ giờ}$, vừa kịp lúc tàu hỏa lăn bánh rời ga. Nếu người đó đi với vận tốc $15\text{ km/h}$ thì đến ga trước giờ tàu chạy bao nhiêu phút?
\choice
{$8\text{ phút}$}
{\True $9\text{ phút}$}
{$10\text{ phút}$}
{$12\text{ phút}$}
\loigiai{
Thời gian người đó đi xe đạp là:
$$7\text{ giờ} - 6\text{ giờ } 15\text{ phút} = 45\text{ phút} = 0{,}75\text{ (giờ)}.$$
Quãng đường từ nhà ra ga tàu là:
$$12 \times 0{,}75 = 9\text{ (km)}.$$
Nếu đi với vận tốc $15\text{ km/h}$ thì thời gian đi là:
$$9 : 15 = 0{,}6\text{ (giờ)} = 36\text{ (phút)}.$$
Người đó đến sớm hơn giờ tàu chạy là:
$$45 - 36 = 9\text{ (phút)}.$$
\textbf{Đáp án của câu hỏi là B}
}
\end{ex}

\begin{ex} %[5D10V62]
Một vận động viên tham gia thi ba môn phối hợp gồm: bơi $1{,}5\text{ km}$ hết $30$ phút, đạp xe $40\text{ km}$ hết $1$ giờ, chạy bộ $10\text{ km}$ hết $45$ phút. Tổng thời gian hoàn thành ba môn phối hợp của vận động viên đó là
\choice
{$2\text{ giờ } 10\text{ phút}$}
{\True $2\text{ giờ } 15\text{ phút}$}
{$2\text{ giờ } 20\text{ phút}$}
{$2\text{ giờ } 30\text{ phút}$}
\loigiai{
Tổng thời gian là:
$$30\text{ phút} + 1\text{ giờ} + 45\text{ phút} = 1\text{ giờ } 75\text{ phút} = 2\text{ giờ } 15\text{ phút}.$$
\textbf{Đáp án của câu hỏi là B}
}
\end{ex}

\begin{ex} %[5D10V62]
Hai thành phố cách nhau $120\text{ km}$. Một ô tô đi từ thành phố A đến B với vận tốc $50\text{ km/h}$, đồng thời một xe máy đi từ B đến A với vận tốc $30\text{ km/h}$. Khi hai xe gặp nhau, vị trí gặp nhau cách B bao nhiêu ki-lô-mét?
\choice
{$40\text{ km}$}
{\True $45\text{ km}$}
{$50\text{ km}$}
{$55\text{ km}$}
\loigiai{
Tổng vận tốc của hai xe là:
$$50 + 30 = 80\text{ (km/h)}.$$
Thời gian hai xe gặp nhau là:
$$120 : 80 = 1{,}5\text{ (giờ)}.$$
Vị trí gặp nhau cách B quãng đường xe máy đi:
$$30 \times 1{,}5 = 45\text{ (km)}.$$
\textbf{Đáp án của câu hỏi là B}
}
\end{ex}

\begin{ex} %[5D10V62]
Kết quả của biểu thức $(4\text{ giờ } 20\text{ phút} - 1\text{ giờ } 45\text{ phút}) \times 2 + 3\text{ giờ } 10\text{ phút}$ là
\choice
{$8\text{ giờ } 10\text{ phút}$}
{\True $8\text{ giờ } 20\text{ phút}$}
{$8\text{ giờ } 30\text{ phút}$}
{$7\text{ giờ } 50\text{ phút}$}
\loigiai{
Trong ngoặc:
$$4\text{ giờ } 20\text{ phút} - 1\text{ giờ } 45\text{ phút} = 3\text{ giờ } 80\text{ phút} - 1\text{ giờ } 45\text{ phút} = 2\text{ giờ } 35\text{ phút}.$$
Nhân với $2$:
$$2\text{ giờ } 35\text{ phút} \times 2 = 4\text{ giờ } 70\text{ phút} = 5\text{ giờ } 10\text{ phút}.$$
Cộng với $3\text{ giờ } 10\text{ phút}$:
$$5\text{ giờ } 10\text{ phút} + 3\text{ giờ } 10\text{ phút} = 8\text{ giờ } 20\text{ phút}.$$
\textbf{Đáp án của câu hỏi là B}
}
\end{ex}

\textbf{Phần 2. Trắc nghiệm đúng sai}
\setcounterpageref{ex}{0}

\begin{ex} %[5D10V62]
Một ô tô và một xe máy cùng xuất phát từ điểm A lúc $7\text{ giờ } 30\text{ phút}$ để đi đến điểm B. Quãng đường AB dài $150\text{ km}$. Vận tốc ô tô là $60\text{ km/h}$, vận tốc xe máy là $40\text{ km/h}$. Xét tính đúng/sai của các nhận định sau:
\choiceTF
{\True Ô tô đến B lúc $10\text{ giờ}$}
{\True Xe máy đến B lúc $11\text{ giờ } 15\text{ phút}$}
{Ô tô đến B trước xe máy $1$ giờ}
{\True Khi ô tô đến B thì xe máy còn cách B một khoảng $50\text{ km}$}
\loigiai{
\begin{itemchoice}
\item Đúng. Thời gian ô tô đi là $150 : 60 = 2{,}5\text{ giờ} = 2\text{ giờ } 30\text{ phút} \Rightarrow 7\text{ giờ } 30\text{ phút} + 2\text{ giờ } 30\text{ phút} = 10\text{ giờ}$.
\item Đúng. Thời gian xe máy đi là $150 : 40 = 3{,}75\text{ giờ} = 3\text{ giờ } 45\text{ phút} \Rightarrow 7\text{ giờ } 30\text{ phút} + 3\text{ giờ } 45\text{ phút} = 11\text{ giờ } 15\text{ phút}$.
\item Sai. Ô tô đến trước xe máy là $3\text{ giờ } 45\text{ phút} - 2\text{ giờ } 30\text{ phút} = 1\text{ giờ } 15\text{ phút}$ (không phải 1 giờ).
\item Đúng. Khi ô tô đến B (sau 2,5 giờ), xe máy đi được $40 \times 2{,}5 = 100\text{ km}$, còn cách B là $150 - 100 = 50\text{ km}$.
\end{itemchoice}
}
\end{ex}

\begin{ex} %[5D10V62]
Cho quãng sông AB dài $40\text{ km}$. Một chiếc ca nô đi xuôi dòng từ A đến B hết $1$ giờ $15$ phút và đi ngược dòng từ B về A hết $2$ giờ. Xét tính đúng/sai của các nhận định sau:
\choiceTF
{\True Vận tốc ca nô khi xuôi dòng là $32\text{ km/h}$}
{\True Vận tốc ca nô khi ngược dòng là $20\text{ km/h}$}
{Vận tốc của dòng nước là $12\text{ km/h}$}
{\True Vận tốc thực của ca nô khi nước yên lặng là $26\text{ km/h}$}
\loigiai{
\begin{itemchoice}
\item Đúng. $1\text{ giờ } 15\text{ phút} = 1{,}25\text{ giờ} \Rightarrow v_{\text{xuôi}} = 40 : 1{,}25 = 32\text{ km/h}$.
\item Đúng. $v_{\text{ngược}} = 40 : 2 = 20\text{ km/h}$.
\item Sai. Vận tốc dòng nước là $(32 - 20) : 2 = 6\text{ km/h}$ (không phải $12\text{ km/h}$).
\item Đúng. Vận tốc thực là $(32 + 20) : 2 = 26\text{ km/h}$.
\end{itemchoice}
}
\end{ex}

\textbf{Phần 3. Câu trả lời ngắn}
\setcounterpageref{ex}{0}

\begin{ex} %[5D10V62]
Hai xe xuất phát cùng một lúc từ hai địa điểm A và B cách nhau $120\text{ km}$ và đi cùng chiều (xe A đuổi theo xe B). Vận tốc xe từ A là $50\text{ km/h}$, vận tốc xe từ B là $30\text{ km/h}$. Hỏi sau bao nhiêu giờ thì xe A đuổi kịp xe B?
\shortans{6}
\loigiai{
Mỗi giờ xe A rút ngắn khoảng cách với xe B là:
$$50 - 30 = 20\text{ (km)}.$$
Thời gian để xe A đuổi kịp xe B là:
$$120 : 20 = 6\text{ (giờ)}.$$
}
\end{ex}

\begin{ex} %[5D10V62]
Một người đi xe máy từ A đến B với vận tốc $40\text{ km/h}$ hết $2$ giờ $15$ phút. Khi trở về từ B đến A, người đó đi với vận tốc $45\text{ km/h}$. Hỏi thời gian người đó đi từ B về A là bao nhiêu phút?
\shortans{120}
\loigiai{
Đổi: $2\text{ giờ } 15\text{ phút} = 2{,}25\text{ giờ}$.\\
Quãng đường AB là:
$$40 \times 2{,}25 = 90\text{ (km)}.$$
Thời gian đi từ B về A là:
$$90 : 45 = 2\text{ (giờ)} = 120\text{ phút}.$$
}
\end{ex}

\begin{ex} %[5D10V62]
Một chiếc xe tải chạy qua một đường hầm dài $700\text{ m}$ với vận tốc $36\text{ km/h}$ mất đúng $1$ phút $15$ giây. Hỏi chiều dài của chiếc xe tải đó là bao nhiêu mét?
\shortans{50}
\loigiai{
Đổi: $36\text{ km/h} = 10\text{ m/s}$.\\
Đổi: $1\text{ phút } 15\text{ giây} = 75\text{ giây}$.\\
Tổng quãng đường xe tải đã đi là:
$$10 \times 75 = 750\text{ (m)}.$$
Chiều dài của chiếc xe tải là:
$$750 - 700 = 50\text{ (m)}.$$
}
\end{ex}

\begin{ex} %[5D10V62]
Hai thành phố cách nhau $175\text{ km}$. Hai ô tô đi ngược chiều nhau cùng lúc xuất phát và gặp nhau sau $2$ giờ $30$ phút. Biết vận tốc xe thứ nhất hơn xe thứ hai là $10\text{ km/h}$. Vận tốc của xe thứ nhất là bao nhiêu ki-lô-mét trên giờ?
\shortans{40}
\loigiai{
Đổi: $2\text{ giờ } 30\text{ phút} = 2{,}5\text{ giờ}$.\\
Tổng vận tốc hai xe là:
$$175 : 2{,}5 = 70\text{ (km/h)}.$$
Vận tốc của xe thứ nhất là:
$$(70 + 10) : 2 = 40\text{ (km/h)}.$$
}
\end{ex}

\begin{ex} %[5D10V62]
Một người hoàn thành $4$ sản phẩm giống nhau hết $5$ giờ $20$ phút. Hỏi người đó hoàn thành $6$ sản phẩm như thế hết bao nhiêu phút?
\shortans{480}
\loigiai{
Đổi: $5\text{ giờ } 20\text{ phút} = 320\text{ phút}$.\\
Thời gian làm 1 sản phẩm là:
$$320 : 4 = 80\text{ (phút)}.$$
Thời gian làm 6 sản phẩm là:
$$80 \times 6 = 480\text{ (phút)}.$$
}
\end{ex}

\begin{ex} %[5D10V62]
Một con chim bồ câu đưa thư bay một quãng đường dài $160\text{ km}$ với vận tốc $64\text{ km/h}$. Thời gian con chim bồ câu bay hết quãng đường đó là bao nhiêu phút?
\shortans{150}
\loigiai{
Thời gian con chim bồ câu bay là:
$$160 : 64 = 2{,}5\text{ (giờ)}.$$
Đổi: $2{,}5\text{ giờ} = 150\text{ phút}$.
}
\end{ex}
